% !TeX root = ../../Book3.tex
\section{多连通区域}
\label{b3:sec:1705}

多连通域的调和函数仍可由边值确定，但共轭的周期多了有限个独立方向。边界分量的调和测度提供了调整这些周期的自然函数。

\subsection{周期}
\label{b3:subsec:170501}

调和函数 \(u\) 的周期为 \(\int_\gamma(-u_y\,\mathrm{d}x+u_x\,\mathrm{d}y)\)。闭形式使周期只依赖曲线同调类。若区域由 \(m+1\) 条互不相交实解析 Jordan 曲线围成，则绕 \(m\) 个内孔的曲线构成独立周期方向；任一闭曲线的周期为这些基本周期的整数线性组合。

\subsection{调和测度}
\label{b3:subsec:170502}

令 \(\omega_j\) 为第 \(j\) 个内边界取一、其余边界取零的调和解。其边界数据在各分量上连续，各边界弧可拉直并取局部线性障碍，所以 Perron 理论保证存在；减去该分量上的常数后作调和反射，还保证解在边界附近光滑。它们满足 \(0<\omega_j<1\)，连同外边界权重 \(\omega_0\) 有 \(\sum_{j=0}^m\omega_j=1\)。因此常数边界数据 \(c_j\) 的解为 \(\sum c_j\omega_j\)。

环域 \(r<|z|<R\) 的内边界权重是 \(\log(R/|z|)/\log(R/r)\)，其共轭周期为 \(-2\pi/\log(R/r)\)。这给出可计算的周期调整系数。

\subsection{调和共轭的障碍}
\label{b3:subsec:170503}

沿正向边界，\(\omega\) 的积分等于调和函数的外法向通量。令 \(A_{ij}=\int_D\nabla\omega_i\cdot\nabla\omega_j\,\mathrm{d}A\)。Green 第一公式把它等于第 \(i\) 个内边界上的 \(\omega_j\) 通量。

矩阵 \(A\) 对称正定：若 \(\sum c_j\omega_j\) 的梯度能量为零，则该组合恒常；外边界值为零使常数零，再看各内边界值知全部 \(c_j=0\)。因此基本周期可以由 \(\omega_j\) 的线性组合任意调整，拓扑障碍有明确的有限维形式。

\subsection{典范表示}
\label{b3:subsec:170504}

\begin{proposition}\label{b3:prop:17-period-decomposition}
在上述有限连通实解析边界有界域中，每个连续到边界的调和函数可写为
\[
u=\operatorname{Re}F+\sum_{j=1}^m c_j\omega_j,
\]
其中 \(F\) 单值全纯，\(c_j\) 唯一；\(F\) 只差纯虚常数。
\end{proposition}
\begin{proof}
用正定通量矩阵解线性方程，选择 \(c_j\) 消除 \(u\) 在全部基本曲线上的周期。剩余调和函数所有闭曲线周期为零，由\cref{b3:thm:15-global-conjugate}得到单值共轭和 \(F\)。若两组系数都可行，其差的周期向量为零，矩阵可逆使差为零；共轭唯一性给最后结论。
\end{proof}


\ProseParagraphBreak

基本周期的方向须统一。沿区域正向边界行进时，区域在左侧，外边界逆时针、内边界顺时针。若切向为 \((x',y')\)，外法向为 \((y',-x')\)，则
\[
(-u_y\,\mathrm{d}x+u_x\,\mathrm{d}y)(\gamma')=\nabla u\cdot n\,\mathrm{d}s.
\]
所以通量矩阵直接给的是顺时针内边界周期；若基本闭曲线统一选逆时针绕孔，则矩阵前须加负号。环域内边界权重的逆时针周期是 \(-2\pi/\log(R/r)\)，而其沿区域正向内边界通量为 \(2\pi/\log(R/r)\)，正好与正定能量相符。

\ProseParagraphBreak

有限维拓扑描述可以通过割线实现。从每个内边界取一条互不相交的简单弧连接到外边界，沿这些弧割开区域后得到单连通域。在割开的域中闭形式有原函数。沿第 \(j\) 条割线两侧的原函数差为绕该孔的周期，且因导数相同而为常数；所有这些差均零时，原函数才能连续拼回原区域。任意闭曲线跨越各条割线的有向次数为整数，其积分等于对应周期的整数和。这同时解释基本周期足够、为何有 \(m\) 个独立方向，以及周期清零后共轭的单值性。

\ProseParagraphBreak

典范表示中的 \(c_j\) 不是边界数据的简单平均。它由整个调和函数的周期决定。环域中取 \(u(z)=\log|z|\)，其周期为 \(2\pi\)，内边界权重周期为 \(-2\pi/\log(R/r)\)，故 \(c=-\log(R/r)\)，余项是常数 \(\log R\)。取 \(u(z)=\operatorname{Re}z\)，周期为零，所以 \(c=0\)，尽管两个边界上的函数值通常不相同。边界权重分解只适合各分量上常数的边值，而周期分解适合一般调和函数；两种线性分解的系数含义不同。

\ProseParagraphBreak

当边界孔数增加，矩阵的大小也增加；如果有无穷多个孔，当前有限维正定论证并未给出一个自动收敛的无穷级数表示。系数、能量和周期的求和需要额外估计。即使孔数有限，若边界粗糙到通常法向导数不存在，也应改用内部周期和弱能量方法，不能继续逐点积分边界法向导数。本节的实解析边界有界范围正是为了使拓扑方向和通量方向能够用同一个有限矩阵联系起来。
